\subsection{Норма сопряженного оператора (в линейном нормированном пространстве)}

\begin{definition}[Сопряженное пространство]
	Пусть \(E\) --- линейное нормированное пространство над полем \(\mathbb{K}\).
	Сопряженным пространством к \(E\) называется пространство всех линейных ограниченных
	функционалов на \(E\):
	\[
	E^*=\mathcal{L}(E,\mathbb{K}).
	\]
\end{definition}

\begin{definition}[Сопряженный оператор]
	Пусть \(E_1,E_2\) --- линейные нормированные пространства и
	\[
	A\in\mathcal{L}(E_1,E_2).
	\]
	Сопряженным к \(A\) называется оператор
	\[
	A^*:E_2^*\to E_1^*,
	\]
	который каждому функционалу \(g\in E_2^*\) сопоставляет функционал
	\[
	A^*g\in E_1^*
	\]
	по правилу
	\[
	(A^*g)(x)=g(Ax),
	\qquad
	x\in E_1.
	\]
	То есть
	\[
	A^*g=g\circ A.
	\]
\end{definition}

\begin{remark}[Следствие Хана--Банаха]
	Для любого элемента \(y\) линейного нормированного пространства \(E\) выполнено
	\[
	\|y\|
	=
	\sup_{\|f\|_{E^*}\leqslant 1}|f(y)|.
	\]
	Именно этот факт будет нужен для доказательства неравенства
	\[
	\|A\|\leqslant \|A^*\|.
	\]
\end{remark}

\begin{theorem}[Норма сопряженного оператора]
	Пусть \(E_1,E_2\) --- линейные нормированные пространства и
	\[
	A\in\mathcal{L}(E_1,E_2).
	\]
	Тогда сопряженный оператор
	\[
	A^*:E_2^*\to E_1^*
	\]
	существует, является линейным ограниченным оператором и
	\[
	\|A^*\|=\|A\|.
	\]
\end{theorem}

\begin{proof}
	Докажем утверждение по смысловым шагам.
	
	\textbf{1. Корректность определения \(A^*g\).}
	
	Пусть
	\[
	g\in E_2^*.
	\]
	Определим отображение
	\[
	A^*g:E_1\to\mathbb{K}
	\]
	по правилу
	\[
	(A^*g)(x)=g(Ax).
	\]
	Проверим, что
	\[
	A^*g\in E_1^*.
	\]
	
	\begin{itemize}
		\item \textbf{Линейность функционала \(A^*g\).}
		
		Пусть
		\[
		x_1,x_2\in E_1,
		\qquad
		\alpha,\beta\in\mathbb{K}.
		\]
		Тогда
		\[
		(A^*g)(\alpha x_1+\beta x_2)
		=
		g(A(\alpha x_1+\beta x_2)).
		\]
		Так как \(A\) линеен,
		\[
		A(\alpha x_1+\beta x_2)
		=
		\alpha Ax_1+\beta Ax_2.
		\]
		Так как \(g\) линеен,
		\[
		g(\alpha Ax_1+\beta Ax_2)
		=
		\alpha g(Ax_1)+\beta g(Ax_2).
		\]
		Следовательно,
		\[
		(A^*g)(\alpha x_1+\beta x_2)
		=
		\alpha(A^*g)(x_1)+\beta(A^*g)(x_2).
		\]
		Значит, \(A^*g\) линеен.
		
		\item \textbf{Ограниченность функционала \(A^*g\).}
		
		Для любого \(x\in E_1\) имеем
		\[
		|(A^*g)(x)|
		=
		|g(Ax)|
		\leqslant
		\|g\|\cdot\|Ax\|.
		\]
		Так как \(A\) ограничен,
		\[
		\|Ax\|\leqslant \|A\|\cdot\|x\|.
		\]
		Поэтому
		\[
		|(A^*g)(x)|
		\leqslant
		\|g\|\cdot\|A\|\cdot\|x\|.
		\]
		Значит,
		\[
		A^*g\in E_1^*.
		\]
	\end{itemize}
	
	\textbf{2. Линейность и ограниченность оператора \(A^*\).}
	
	Теперь рассматриваем \(A^*\) уже как оператор
	\[
	A^*:E_2^*\to E_1^*.
	\]
	
	\begin{itemize}
		\item \textbf{Линейность оператора \(A^*\).}
		
		Пусть
		\[
		g_1,g_2\in E_2^*,
		\qquad
		\alpha,\beta\in\mathbb{K}.
		\]
		Тогда для любого \(x\in E_1\)
		\[
		(A^*(\alpha g_1+\beta g_2))(x)
		=
		(\alpha g_1+\beta g_2)(Ax).
		\]
		Отсюда
		\[
		(A^*(\alpha g_1+\beta g_2))(x)
		=
		\alpha g_1(Ax)+\beta g_2(Ax)
		=
		\alpha(A^*g_1)(x)+\beta(A^*g_2)(x).
		\]
		Так как равенство верно для всех \(x\in E_1\), получаем
		\[
		A^*(\alpha g_1+\beta g_2)
		=
		\alpha A^*g_1+\beta A^*g_2.
		\]
		Значит, \(A^*\) линеен.
		
		\item \textbf{Ограниченность оператора \(A^*\).}
		
		Из оценки первого пункта
		\[
		|(A^*g)(x)|
		\leqslant
		\|g\|\cdot\|A\|\cdot\|x\|
		\]
		получаем
		\[
		\|A^*g\|
		\leqslant
		\|A\|\cdot\|g\|.
		\]
		Следовательно,
		\[
		A^*\in\mathcal{L}(E_2^*,E_1^*)
		\]
		и
		\[
		\|A^*\|\leqslant\|A\|.
		\]
	\end{itemize}
	
	\textbf{3. Обратное неравенство \(\|A\|\leqslant\|A^*\|\).}
	
	Возьмём произвольный элемент
	\[
	x\in E_1.
	\]
	По следствию Хана--Банаха для элемента \(Ax\in E_2\)
	\[
	\|Ax\|
	=
	\sup_{\|g\|_{E_2^*}\leqslant1}|g(Ax)|.
	\]
	По определению сопряженного оператора
	\[
	g(Ax)=(A^*g)(x).
	\]
	Поэтому
	\[
	\|Ax\|
	=
	\sup_{\|g\|_{E_2^*}\leqslant1}|(A^*g)(x)|.
	\]
	Для каждого \(g\in E_2^*\), такого что \(\|g\|\leqslant1\), имеем
	\[
	|(A^*g)(x)|
	\leqslant
	\|A^*g\|\cdot\|x\|.
	\]
	Но
	\[
	\|A^*g\|
	\leqslant
	\|A^*\|\cdot\|g\|
	\leqslant
	\|A^*\|.
	\]
	Значит,
	\[
	|(A^*g)(x)|
	\leqslant
	\|A^*\|\cdot\|x\|.
	\]
	Переходя к супремуму по всем \(\|g\|\leqslant1\), получаем
	\[
	\|Ax\|
	\leqslant
	\|A^*\|\cdot\|x\|.
	\]
	Так как это верно для любого \(x\in E_1\), то
	\[
	\|A\|\leqslant\|A^*\|.
	\]
	
	\textbf{4. Равенство норм.}
	
	Мы доказали две оценки:
	\[
	\|A^*\|\leqslant\|A\|
	\]
	и
	\[
	\|A\|\leqslant\|A^*\|.
	\]
	Следовательно,
	\[
	\|A^*\|=\|A\|.
	\]
\end{proof}
